\subsection{الگوریتم بهینه‌سازی تقریبی خط‌مشی (\en{Proximal Policy Optimization - PPO})}

الگوریتم \model{PPO} که توسط جان شولمن (\en{John Schulman}) و همکارانش در سال ۲۰۱۷ در \en{OpenAI} معرفی شد \cite{schulman2017proximal}، یکی از برجسته‌ترین، پایدارترین و پرکاربردترین الگوریتم‌های حوزه یادگیری تقویتی عمیق (\en{Deep Reinforcement Learning}) به شمار می‌رود. این الگوریتم با حل چالش‌های بنیادین بهینه‌سازی در روش‌های پیشین، تعادلی بهینه میان سادگی پیاده‌سازی، کارایی محاسباتی، بهره‌وری نمونه (\en{Sample Efficiency}) و پایداری همگرایی برقرار ساخته است.

\subsubsection*{تاریخچه، انگیزه و ضرورت پیدایش الگوریتم}

پیش از ابداع \model{PPO}، روش‌های یادگیری تقویتی عمیق به سه دسته عمده با محدودیت‌های ساختاری تقسیم می‌شدند:
\begin{itemize}
    \item \textbf{روش‌های مبتنی بر ارزش (\en{Value-Based Methods}):} شاخص‌ترین الگوریتم این دسته، \model{DQN} \cite{mnih2015human} است. این روش‌ها گرچه در فضاهای عمل گسسته و محیط بازی‌های آتاری \cite{bellemare2015arcade} موفق ظاهر شدند، اما امکان تعمیم کارآمد به فضاهای عمل پیوسته (\en{Continuous Action Spaces}) با ابعاد بالا در رباتیک \cite{brockman2016openai, duan2016benchmarking} را نداشتند، چرا که یافتن $\arg\max_a Q(s, a)$ در هر گام نیازمند بهینه‌سازی غیرخطی زمان‌بر است. همچنین این روش‌ها از ناپایداری‌های ناشی از تقریب تابع و یادگیری خارج از خط‌مشی رنج می‌برند.
    \item \textbf{روش‌های گرادیان خط‌مشی استاندارد (\en{Vanilla Policy Gradient}):} روش‌هایی نظیر \en{REINFORCE} \cite{williams1992simple} و \model{A2C/A3C} \cite{mnih2016asynchronous} مستقیماً پارامترهای خط‌مشی را بهینه‌سازی می‌کنند. با این حال، به دلیل وابستگی گرادیان به توزیع داده‌های جمع‌آوری‌شده، مجاز به اجرای تنها یک گام گرادیان به ازای هر نمونه داده بودند. انجام چندین گام بهینه‌سازی (\en{Multiple Epochs}) روی داده‌های یکسان، منجر به جابجایی شدید خط‌مشی جدید نسبت به خط‌مشی قدیم و در نتیجه سقوط فاجعه‌بار کارایی (\en{Catastrophic Policy Collapse}) می‌شد؛ امری که باعث ناکارآمدی شدید در تعداد نمونه‌ها (\en{Sample Inefficiency}) بود.
    \item \textbf{روش‌های بهینه‌سازی در ناحیه اطمینان (\en{Trust Region Methods}):} کاکاده و لنگفورد \cite{kakade2002approximately} و سپس شولمن و همکاران با معرفی الگوریتم \model{TRPO} \cite{schulman2015trust} نشان دادند که با مقید ساختن تغییرات خط‌مشی از طریق واگرایی کولبک-لایبلر (\en{KL-Divergence})، می‌توان بهبود یکنواخت کارایی (\en{Monotonic Improvement}) را تضمین کرد. با این وجود، \model{TRPO} به محاسبات مرتبه دوم ماتریس اطلاعات فیشر (\en{Fisher Information Matrix})، ماتریس هسین و حل با الگوریتم گرادیان مزدوج (\en{Conjugate Gradient}) نیازمند بود. این امر نه تنها پیچیدگی محاسباتی بالایی ایجاد می‌کرد، بلکه با معماری‌های دارای لایه‌های مشترک میان خط‌مشی و ارزش و تکنیک‌های نویزدار مانند \en{Dropout} ناسازگار بود.
\end{itemize}

هدف بنیادین \model{PPO}، دستیابی به ثبات، قابلیت اطمینان و نرخ همگرایی بالای \model{TRPO}، صرفاً با تکیه بر بهینه‌سازهای مشتق اول (\en{First-Order Optimizers}) نظیر بهینه‌ساز \en{Adam} \cite{kingma2014adam} و با سادگی پیاده‌سازی بالا است.

\subsubsection*{مبانی نظری و اشتقاق‌های کامل ریاضی}

یک فرآیند تصمیم‌گیری مارکوف (\en{MDP}) با چندتایی متناهی یا نامتناهی $\mathcal{M} = (\mathcal{S}, \mathcal{A}, \mathcal{P}, r, \gamma, \rho_0)$ مدل‌سازی می‌شود که در آن $\mathcal{S}$ فضای حالات، $\mathcal{A}$ فضای اعمال، $\mathcal{P}(s'|s, a)$ تابع احتمال انتقال، $r(s, a)$ تابع پاداش، $\gamma \in [0, 1)$ ضریب تخفیف پاداش و $\rho_0(s_0)$ توزیع حالت اولیه است.

توزیع یک مسیر (\en{Trajectory}) $\tau = (s_0, a_0, s_1, a_1, \dots, s_T)$ تحت خط‌مشی تصادفی پارامتریزه شده $\pi_\theta(a|s)$ برابر است با:
\begin{equation}
	P(\tau; \theta) = \rho_0(s_0) \prod_{t=0}^{T-1} \pi_\theta(a_t | s_t) \mathcal{P}(s_{t+1} | s_t, a_t)
\end{equation}

هدف بهینه‌سازی، بیشینه‌سازی امید ریاضی بازده کل است:
\begin{equation}
	J(\theta) = \mathbb{E}_{\tau \sim P(\cdot; \theta)} [R(\tau)] = \int_{\tau} P(\tau; \theta) R(\tau) \, d\tau
\end{equation}
که در آن $R(\tau) = \sum_{t=0}^T \gamma^t r(s_t, a_t)$ است. با اعمال عملگر گرادیان و استفاده از قاعده مشتق لگاریتمی ($\nabla_\theta P = P \nabla_\theta \log P$):
\begin{align}
	\nabla_\theta J(\theta) &= \nabla_\theta \int P(\tau; \theta) R(\tau) \, d\tau = \int \nabla_\theta P(\tau; \theta) R(\tau) \, d\tau \nonumber \\
	&= \int P(\tau; \theta) \nabla_\theta \log P(\tau; \theta) R(\tau) \, d\tau = \mathbb{E}_{\tau \sim \pi_\theta} \left[ \nabla_\theta \log P(\tau; \theta) R(\tau) \right]
\end{align}
از آنجا که جملات انتقال وضعیت محیط $\mathcal{P}$ به پارامتر $\theta$ وابسته نیستند، داریم $\nabla_\theta \log P(\tau; \theta) = \sum_{t=0}^T \nabla_\theta \log \pi_\theta(a_t | s_t)$. با بهره‌گیری از اصل علیت و کسر خط مبنای مستقل از عمل $b(s_t) = V(s_t)$ جهت کاهش واریانس، قضیه گرادیان خط‌مشی حاصل می‌شود:
\begin{equation}
	\hat{g} = \hat{\mathbb{E}}_t \left[ \nabla_\theta \log \pi_\theta(a_t | s_t) \hat{A}_t \right]
\end{equation}
که در آن $\hat{A}_t = Q(s_t, a_t) - V(s_t)$ برآوردگر تابع مزیت (\en{Advantage Function}) است.

با استفاده از مفهوم نمونه‌برداری ترجیحی (\en{Importance Sampling})، هدف خط‌مشی برای داده‌های جمع‌آوری‌شده توسط $\pi_{\theta_{\text{old}}}$ به صورت تابع هدف جایگزین تکرار خط‌مشی محافظه‌کارانه (\en{CPI}) \cite{kakade2002approximately} تعریف می‌شود:
\begin{equation}
	L^{CPI}(\theta) = \hat{\mathbb{E}}_t \left[ \frac{\pi_\theta(a_t | s_t)}{\pi_{\theta_{\text{old}}}(a_t | s_t)} \hat{A}_t \right] = \hat{\mathbb{E}}_t \left[ r_t(\theta) \hat{A}_t \right]
\end{equation}
که در آن نسبت احتمالات $r_t(\theta) = \frac{\pi_\theta(a_t | s_t)}{\pi_{\theta_{\text{old}}}(a_t | s_t)}$ بوده و $r_t(\theta_{\text{old}}) = 1$ است.

\subsubsection*{توابع هدف در الگوریتم PPO}

الگوریتم \model{PPO} دو فرمول‌بندی برای محدود ساختن تغییرات گام خط‌مشی ارائه می‌دهد:

\textbf{۱. تابع هدف با برش جایگزین (\en{Clipped Surrogate Objective}):}
رویکرد اصلی و پیشنهادی مقاله، استفاده از تابع هدف زیر است:
\begin{equation}
	L^{\text{CLIP}}(\theta) = \hat{\mathbb{E}}_t \left[ \min \left( r_t(\theta) \hat{A}_t, \; \text{clip}(r_t(\theta), 1-\epsilon, 1+\epsilon) \hat{A}_t \right) \right]
\end{equation}
که در آن $\epsilon$ یک ابرپارامتر (معمولاً $\epsilon = 0.2$) است. این تابع هدف یک کران پایین بدبینانه (\en{Pessimistic Lower Bound}) را تشکیل می‌دهد:
\begin{itemize}
	\item \textbf{در شرایط مزیت مثبت ($\hat{A}_t > 0$):} اگر عمل انجام‌شده بهتر از حد انتظار باشد، افزایش احتمال آن مفید است؛ اما به محض اینکه نسبت احتمال از $1+\epsilon$ فراتر رود، گرادیان صفر شده و از تغییرات مفرط جلوگیری می‌کند.
	\item \textbf{در شرایط مزیت منفی ($\hat{A}_t < 0$):} اگر عمل بدتر از حد انتظار باشد، احتمال آن باید کاهش یابد؛ اگر نسبت احتمال از $1-\epsilon$ کمتر شود، مقدار بریده شده و گرادیان صفر می‌شود. اما اگر خط‌مشی اشتباهاً احتمال این عمل بد را افزایش دهد ($r_t \gg 1$)، مقدار بدون برش انتخاب شده و جریمه کامل به خط‌مشی اعمال می‌شود.
\end{itemize}

\textbf{۲. تابع هدف با ضریب جریمه تطبیقی (\en{Adaptive KL Penalty}):}
به عنوان رویکرد جایگزین، جریمه واگرایی \en{KL} مستقیماً در تابع هدف لحاظ شده و ضریب $\beta$ به صورت دینامیک به‌روزرسانی می‌شود:
\begin{equation}
	L^{\text{KLPEN}}(\theta) = \hat{\mathbb{E}}_t \left[ r_t(\theta)\hat{A}_t - \beta \, D_{\text{KL}}\left(\pi_{\theta_{\text{old}}}(\cdot|s_t) \parallel \pi_\theta(\cdot|s_t)\right) \right]
\end{equation}
در هر تکرار، پس از محاسبه $d = \hat{\mathbb{E}}_t [D_{\text{KL}}(\pi_{\theta_{\text{old}}} \parallel \pi_\theta)]$، ضریب جریمه به این صورت تطبیق می‌یابد:
\begin{equation}
	\beta \leftarrow \begin{cases} 
		\beta / 2 & \text{if } d < d_{\text{targ}} / 1.5 \\ 
		\beta \times 2 & \text{if } d > d_{\text{targ}} \times 1.5 
	\end{cases}
\end{equation}

\subsubsection*{برآوردگر مزیت تعمیم‌یافته (\en{GAE}) و تابع هزینه یکپارچه}

برای کاهش واریانس، از برآوردگر مزیت تعمیم‌یافته (\en{GAE}) منقطع \cite{schulman2015high} استفاده می‌شود:
\begin{equation}
	\hat{A}_t = \sum_{l=0}^{T-t-1} (\gamma \lambda)^l \delta_{t+l}^V \quad \text{where} \quad \delta_t^V = r_t + \gamma V(s_{t+1}) - V(s_t)
\end{equation}
برای شبکه‌های دارای پارامترهای مشترک میان خط‌مشی و تابع ارزش، تابع هزینه سراسری بهینه‌سازی زیر بیشینه می‌گردد:
\begin{equation}
	L_t^{\text{CLIP+VF+S}}(\theta) = \hat{\mathbb{E}}_t \left[ L_t^{\text{CLIP}}(\theta) - c_1 \left( V_\theta(s_t) - V_t^{\text{targ}} \right)^2 + c_2 S[\pi_\theta](s_t) \right]
\end{equation}
که در آن $S[\pi_\theta]$ آنتروپی شانون برای تشویق به اکتشاف (\en{Exploration}) و $c_1, c_2$ ضرایب وزنی هستند.

\subsubsection*{تحلیل‌های تئوری اطلاعات و آمار}

بر اساس بسط تیلور مرتبه دوم واگرایی \en{KL} حول نقطه $\theta_{\text{old}}$ داریم:
\begin{equation}
	D_{\text{KL}}(\pi_{\theta_{\text{old}}} \parallel \pi_\theta) \approx \frac{1}{2} (\theta - \theta_{\text{old}})^T F (\theta - \theta_{\text{old}})
\end{equation}
که در آن $F = \mathbb{E} [\nabla_\theta \log \pi_\theta \nabla_\theta \log \pi_\theta^T]$ ماتریس اطلاعات فیشر است. تابع برش در \model{PPO} تضمین می‌کند که فاصله روی منیفلد ریمانی توزیع‌ها بدون نیاز به تشکیل و معکوس کردن صریح ماتریس $F$ در محدوده‌ای امن محبوس بماند. همچنین با محدود کردن $r_t(\theta) \in [1-\epsilon, 1+\epsilon]$، واریانس بازوزن‌دهی نمونه‌برداری ترجیحی ($\operatorname{Var}_{Q}[P/Q \cdot f]$) کنترل شده و از انفجار واریانس در بهینه‌سازی چند دوره‌ای جلوگیری می‌شود.

\subsubsection*{دیاگرام معماری و جریان داده‌های الگوریتم}

شکل~\ref{fig:ppo_architecture}، ساختار تعامل با محیط، محاسبه مزیت، تشکیل توابع هدف و چرخه به‌روزرسانی چند دوره‌ای در الگوریتم \model{PPO} را نشان می‌دهد.

\begin{figure}[htbp]
\centering
\begin{latin}
\begin{tikzpicture}[
	node distance=1.4cm and 1.8cm,
	box/.style={rectangle, rounded corners=6pt, draw=blue!70!black, fill=blue!5, thick, inner sep=8pt, align=center, font=\small\sffamily, drop shadow},
	envbox/.style={rectangle, rounded corners=6pt, draw=teal!70!black, fill=teal!10, thick, inner sep=8pt, align=center, font=\small\sffamily, drop shadow},
	lossbox/.style={rectangle, rounded corners=6pt, draw=purple!70!black, fill=purple!5, thick, inner sep=8pt, align=center, font=\small\sffamily, drop shadow},
	optbox/.style={rectangle, rounded corners=6pt, draw=orange!80!black, fill=orange!10, thick, inner sep=8pt, align=center, font=\small\sffamily, drop shadow},
	arrow/.style={-{Stealth[scale=1.1]}, thick, draw=blue!80!black},
	dasharrow/.style={-{Stealth[scale=1.1]}, thick, dashed, draw=red!70!black}
]

	% Environment & Actors
	\node[envbox] (env) {\textbf{Parallel Environments}\\ $N$ Actors running $\pi_{\theta_{\text{old}}}$ for $T$ timesteps};
	
	% Memory Buffer
	\node[box, below=of env] (buffer) {\textbf{Trajectory Buffer}\\ Store states $s_t$, actions $a_t$, rewards $r_t$,\\ log-probs $\log\pi_{\theta_{\text{old}}}(a_t|s_t)$, values $V(s_t)$};
	
	% GAE Estimation
	\node[box, below=of buffer] (gae) {\textbf{GAE Advantage Engine}\\ $\delta_t^V = r_t + \gamma V(s_{t+1}) - V(s_t)$\\ $\hat{A}_t = \sum_{l=0}^{T-t-1} (\gamma\lambda)^l \delta_{t+l}^V$, \quad $V_t^{\text{targ}} = \hat{A}_t + V(s_t)$};

	% Surrogate Loss Construction
	\node[lossbox, below left=1.5cm and -0.5cm of gae] (lclip) {\textbf{Clipped Policy Loss}\\ $r_t(\theta) = \frac{\pi_\theta(a_t|s_t)}{\pi_{\theta_{\text{old}}}(a_t|s_t)}$\\ $L^{\text{CLIP}} = \hat{\mathbb{E}}[\min(r_t\hat{A}_t, \text{clip}(r_t, 1\pm\epsilon)\hat{A}_t)]$};
	
	\node[lossbox, below right=1.5cm and -0.5cm of gae] (lvf) {\textbf{Value \& Entropy Loss}\\ $L^{VF} = (V_\theta(s_t) - V_t^{\text{targ}})^2$\\ $S[\pi_\theta] = -\sum \pi_\theta \log \pi_\theta$};

	% Total Loss
	\node[lossbox, below=3.8cm of gae] (ltot) {\textbf{Total Surrogate Objective}\\ $L^{\text{Total}}(\theta) = L^{\text{CLIP}}(\theta) - c_1 L^{VF}(\theta) + c_2 S[\pi_\theta]$};

	% Optimizer
	\node[optbox, below=of ltot] (adam) {\textbf{Minibatch SGD / Adam Optimizer}\\ Perform $K$ Epochs on minibatches $M \le NT$\\ Update $\theta \leftarrow \theta + \alpha \nabla_\theta L^{\text{Total}}(\theta)$};

	% Reference Update
	\node[envbox, left=2.2cm of gae] (sync) {\textbf{Policy Synchronization}\\ $\theta_{\text{old}} \leftarrow \theta$};

	% Connections
	\draw[arrow] (env) -- node[right, font=\footnotesize] {Transitions} (buffer);
	\draw[arrow] (buffer) -- node[right, font=\footnotesize] {Data} (gae);
	\draw[arrow] (gae) -| (lclip);
	\draw[arrow] (gae) -| (lvf);
	\draw[arrow] (lclip) |- (ltot);
	\draw[arrow] (lvf) |- (ltot);
	\draw[arrow] (ltot) -- node[right, font=\footnotesize] {$\nabla_\theta L^{\text{Total}}$} (adam);
	
	% Loop back
	\draw[dasharrow] (adam) -| (sync);
	\draw[dasharrow] (sync) |- (env);

\end{tikzpicture}
\end{latin}
\caption{دیاگرام جامع جریان داده‌ها، محاسبه مزیت و بهینه‌سازی چند دوره‌ای در الگوریتم \model{PPO}.}
\label{fig:ppo_architecture}
\end{figure}

\subsubsection*{شبه‌کد الگوریتم PPO}

شبه‌کد استاندارد ساختار \en{Actor-Critic} الگوریتم \model{PPO} در الگوریتم~\ref{alg:ppo_actor_critic} آورده شده است:

\begin{latin}
\begin{algorithm}[H]
\caption{Proximal Policy Optimization (PPO), Actor-Critic Style}
\label{alg:ppo_actor_critic}
\begin{algorithmic}[1]
\State \textbf{Input:} Initial policy parameters $\theta_0$, initial value function parameters $\phi_0$, clipping threshold $\epsilon$, GAE parameters $\gamma, \lambda$, epoch count $K$, minibatch size $M$, coefficients $c_1, c_2$.
\For{$\text{iteration} = 1, 2, \dots$}
    \For{$\text{actor} = 1, \dots, N$}
        \State Run policy $\pi_{\theta_{\text{old}}}$ in environment for $T$ timesteps.
        \State Compute predictions $V_\phi(s_t)$ and log-probabilities $\log \pi_{\theta_{\text{old}}}(a_t | s_t)$.
        \State Compute TD residuals: $\delta_t^V = r_t + \gamma V_\phi(s_{t+1}) - V_\phi(s_t)$.
        \State Compute GAE estimates: $\hat{A}_t = \sum_{l=0}^{T-t-1} (\gamma \lambda)^l \delta_{t+l}^V$.
        \State Compute state-value targets: $V_t^{\text{targ}} = \hat{A}_t + V_\phi(s_t)$.
    \EndFor
    \State Collect all $NT$ transitions: $\mathcal{D} = \{(s_t, a_t, r_t, \log\pi_{\theta_{\text{old}}}(a_t|s_t), \hat{A}_t, V_t^{\text{targ}})\}$.
    \State Normalize advantage values: $\hat{A}_t \leftarrow \frac{\hat{A}_t - \mu(\hat{A})}{\sigma(\hat{A}) + 10^{-8}}$.
    \For{$\text{epoch} = 1, \dots, K$}
        \State Shuffle dataset $\mathcal{D}$ and partition into minibatches of size $M$.
        \For{each minibatch $\mathcal{B} \subset \mathcal{D}$}
            \State Compute probability ratio: $r_t(\theta) = \exp\left(\log \pi_\theta(a_t|s_t) - \log \pi_{\theta_{\text{old}}}(a_t|s_t)\right)$.
            \State Compute policy surrogate loss: 
            \State \quad $L^{\text{CLIP}}(\theta) = \frac{1}{|\mathcal{B}|} \sum_{t \in \mathcal{B}} \min\left( r_t(\theta)\hat{A}_t, \, \text{clip}(r_t(\theta), 1-\epsilon, 1+\epsilon)\hat{A}_t \right)$.
            \State Compute value function squared error loss:
            \State \quad $L^{\text{VF}}(\phi) = \frac{1}{|\mathcal{B}|} \sum_{t \in \mathcal{B}} \left( V_\phi(s_t) - V_t^{\text{targ}} \right)^2$.
            \State Compute entropy bonus: $S[\pi_\theta] = \frac{1}{|\mathcal{B}|} \sum_{t \in \mathcal{B}} \mathbb{H}(\pi_\theta(\cdot | s_t))$.
            \State Maximize total objective using Adam:
            \State \quad $\theta, \phi \leftarrow \text{Adam}\left( \nabla_{\theta, \phi} \left[ L^{\text{CLIP}}(\theta) - c_1 L^{\text{VF}}(\phi) + c_2 S[\pi_\theta] \right] \right)$.
        \EndFor
    \EndFor
    \State $\theta_{\text{old}} \leftarrow \theta$
\EndFor
\end{algorithmic}
\end{algorithm}
\end{latin}

\subsubsection*{نتایج تجربی و ارزیابی بنچمارک‌ها}

کارایی \model{PPO} در مقایسه با انواع روش‌های بهینه‌سازی در جدول~\ref{tab:mujoco_ablation} و روی بازی‌های آتاری در جدول~\ref{tab:atari_comparison} خلاصه شده است:

\begin{table}[htbp]
\centering
\caption{مقایسه عملکرد انواع توابع هدف جایگزین در ۷ محیط پیوسته \en{MuJoCo} (برگرفته از مقاله مرجع \cite{schulman2017proximal}).}
\label{tab:mujoco_ablation}
\begin{tabular}{lc}
\toprule
\textbf{فرمول‌بندی الگوریتم / تابع هدف} & \textbf{میانگین امتیاز نرمال‌شده} \\
\midrule
بدون برش و بدون جریمه (\en{No clipping or penalty}) & -0.39 \\
برش با $\epsilon = 0.1$ (\en{Clipping}) & 0.76 \\
\textbf{برش با $\epsilon = 0.2$ (\en{Clipping})} & \textbf{0.82} \\
برش با $\epsilon = 0.3$ (\en{Clipping}) & 0.70 \\
جریمه تطبیقی با $d_{\text{targ}} = 0.003$ (\en{Adaptive KL}) & 0.68 \\
جریمه تطبیقی با $d_{\text{targ}} = 0.01$ (\en{Adaptive KL}) & 0.74 \\
جریمه تطبیقی با $d_{\text{targ}} = 0.03$ (\en{Adaptive KL}) & 0.71 \\
جریمه ثابت با $\beta = 0.3$ (\en{Fixed KL}) & 0.62 \\
جریمه ثابت با $\beta = 1.0$ (\en{Fixed KL}) & 0.71 \\
جریمه ثابت با $\beta = 3.0$ (\en{Fixed KL}) & 0.72 \\
جریمه ثابت با $\beta = 10.0$ (\en{Fixed KL}) & 0.69 \\
\bottomrule
\end{tabular}
\end{table}

\begin{table}[htbp]
\centering
\caption{تعداد بازی‌های برده‌شده توسط هر الگوریتم در بنچمارک ۴۹ تایی آتاری (\en{ALE}) \cite{schulman2017proximal}.}
\label{tab:atari_comparison}
\begin{tabular}{lcccc}
\toprule
\textbf{معیار ارزیابی} & \textbf{\model{A2C}} & \textbf{\model{ACER}} & \textbf{\model{PPO}} & \textbf{تساوی} \\
\midrule
(۱) میانگین پاداش در کل طول آموزش & 1 & 18 & \textbf{30} & 0 \\
(۲) میانگین پاداش در ۱۰۰ اپیزود پایانی & 1 & \textbf{28} & 19 & 1 \\
\bottomrule
\end{tabular}
\end{table}

\subsubsection*{ابرپارامترهای استاندارد گزارش‌شده در مقاله}

جداول~\ref{tab:params_mujoco} و \ref{tab:params_atari} مقادیر دقیق ابرپارامترهای بهینه‌شده برای محیط‌های مختلف را نمایش می‌دهند:

\begin{table}[htbp]
\centering
\caption{ابرپارامترهای \model{PPO} در بنچمارک رباتیک پیوسته \en{MuJoCo} (۱ میلیون گام).}
\label{tab:params_mujoco}
\begin{tabular}{lc}
\toprule
\textbf{ابرپارامتر} & \textbf{مقدار} \\
\midrule
افق زمانی داده‌برداری ($T$) & 2048 \\
نرخ یادگیری بهینه‌ساز \en{Adam} & $3 \times 10^{-4}$ \\
تعداد دور بهینه‌سازی در هر آپدیت ($K$) & 10 \\
اندازه دسته کوچک ($M$) & 64 \\
ضریب تخفیف پاداش ($\gamma$) & 0.99 \\
پارامتر برآوردگر مزیت ($\lambda$) & 0.95 \\
\bottomrule
\end{tabular}
\end{table}

\begin{table}[htbp]
\centering
\caption{ابرپارامترهای \model{PPO} در بنچمارک بازی‌های آتاری (\en{Atari}).}
\label{tab:params_atari}
\begin{tabular}{lc}
\toprule
\textbf{ابرپارامتر} & \textbf{مقدار} \\
\midrule
افق زمانی داده‌برداری ($T$) & 128 \\
نرخ یادگیری \en{Adam} & $2.5 \times 10^{-4} \times \alpha$ \\
تعداد عامل‌های موازی ($N$) & 8 \\
تعداد دورها ($K$) & 3 \\
اندازه دسته کوچک ($M$) & $32 \times 8 = 256$ \\
ضریب تخفیف پاداش ($\gamma$) & 0.99 \\
پارامتر $\lambda$ در \en{GAE} & 0.95 \\
آستانه برش ($\epsilon$) & $0.1 \times \alpha$ \\
ضریب خطای تابع ارزش ($c_1$) & 1.0 \\
ضریب پاداش آنتروپی ($c_2$) & 0.01 \\
\bottomrule
\end{tabular}
\end{table}

\subsubsection*{جمع‌بندی}
الگوریتم \model{PPO} به واسطه تابع برش بدبینانه $L^{\text{CLIP}}$، بدون نیاز به محاسبات سنگین ماتریس فیشر و با استفاده مستقیم از گرادیان مرتبه اول \en{Adam}، رفتار ناحیه اطمینان را شبیه‌سازی کرده و امروزه به عنوان پایه‌ای‌ترین و مطمئن‌ترین الگوریتم یادگیری تقویتی در زمینه‌های مختلف از کنترل حرکت ربات‌های سه‌بعدی \cite{heess2017emergence} تا هم‌ترازسازی مدل‌های زبانی بزرگ از طریق بازخورد انسانی (\en{RLHF}) شناخته می‌شود.